\section{A deck from Grothendieck's notes: pseudo-fields}\label{sec:pseudofields}

The precedents of \S\ref{sec:precedents} show Grothendieck reading invariants of a topos on two presentations. His notes on universal algebra also contain a deck that the registry did not have: the topos of \q{pseudo-fields}, with three presentations of it.

\subsection{The deck}

In folder 161-2 \cite{grothendieck-161-2}, pp.~38--43, Grothendieck lists properties of rings and, for each, a subtopos of the topos $X_{\mathrm{ann}}\simeq\widehat R$ classifying rings, where $R$ is the category of affine schemes of finite type over $\mathbb Z$. Item~13 (p.~42) is the \emph{compact reduced ring}, which he names \emph{pseudo-field}: for every $f$ there is a unique idempotent $e$ with $ef$ invertible in $eA$ and $(1-e)f=0$, so that $A\simeq A_f\times A/fA$. These are the commutative von Neumann regular rings. He writes three presentations of the subtopos they define:
\[
X_{\mathrm{ps.\,corps}}\;\simeq\;\widetilde R_{\mathrm{ps.\,corps}}\;\simeq\;\widehat{R_{\mathrm{cons}}},
\]
sheaves on $R$ for the topology in which a family covers when one of its members has a section fibre by fibre, and presheaves on the category $R_{\mathrm{cons}}$ obtained from $R$ by inverting the arrows that are bijective with trivial residue field extensions (pp.~41--42).

\subsection{Checking the presentation by presheaves}

Since the idempotent $e$ is unique, the theory of pseudo-fields is cartesian, hence of presheaf type: its classifying topos is the topos of $\Set$-valued functors on its finitely presentable models. These are the absolutely flat hulls $T(A)$ of the finitely presented rings $A$, the left adjoint $T$ of the inclusion of von Neumann regular rings into rings preserving finite presentability. By Olivier \cite{olivier}, $\operatorname{Spec}T(A)$ is $\operatorname{Spec}A$ with its constructible topology, and the local rings of $T(A)$ are the residue fields of $A$. By Pierce \cite{pierce}, a von Neumann regular ring is the ring of global sections of the sheaf of its local rings over its spectrum, a Boolean space. Hence for $u:A\to B$ of finite presentation, $T(u)$ is an isomorphism iff $\operatorname{Spec}u$ is bijective and every residue field extension is trivial: a continuous bijection of Boolean spaces is a homeomorphism, and a map of sheaves that is an isomorphism on stalks is an isomorphism. The arrows that $T$ inverts are therefore exactly those that Grothendieck inverts, which is the presentation $\widehat{R_{\mathrm{cons}}}$ read through the finitely presentable models. We have not checked the first presentation, by the topology $\mathrm{ps.\,corps}$, and nothing in this section is formalised.

\subsection{The arches on this deck}

\begin{itemize}
\item \emph{De Morgan:} yes; the finitely presentable models of a cartesian theory have pushouts, so they amalgamate (Bridge~\ref{br:dm}).
\item \emph{Two-valued:} no; the sentence $1=0$ defines a subterminal object that is neither $0$ nor $1$.
\item \emph{Connected:} yes; $T(\mathbb Z)$ is an initial finitely presentable model, so the category of these models is connected.
\item \emph{Boolean:} no; the map $T(\mathbb Z)\to T(\mathbb F_2)$ is not invertible.
\item \emph{Essential points:} the finitely presentable pseudo-fields $T(A)$, up to isomorphism.
\item \emph{Subtopoi:} the quotient theories of the theory of pseudo-fields. Grothendieck's next items are such quotients: fields, which he identifies with the integral pseudo-fields and with the local ones (item~12, p.~41, the last equality marked \q{à prouver}, and true, since a local von Neumann regular ring is a field); separably and algebraically closed fields (items~14--15); rings killed by $n$ and rings in which $n$ is invertible (item~16, p.~43), which he marks as open subtopoi.
\end{itemize}

\subsection{An answer to a question in the margin}

For item~12, the \emph{compact} rings (every $f$ invertible on an idempotent $e$ and nilpotent on $1-e$, p.~42), Grothendieck notes in the margin that the topology is not coherent and asks whether the topos has enough points. On p.~41 he also warns that the subtopoi of his list are determined by their points, but that this is not clear for their intersections, so that checking an inclusion between intersections on set-theoretic structures may not suffice. Both doubts are settled by a theorem of Makkai and Reyes \cite[Thm~6.2.4]{makkai-reyes} (see also \cite[Cor.~3.3.10]{diliberti-rogers}): a topos of sheaves on a countable category with pullbacks, for a topology generated by a countable family of sieves, has enough points. The category $R$ is essentially countable and has pullbacks, and each of Grothendieck's topologies is generated by countably many covering families, as is any intersection of countably many of his subtopoi. All these topoi have enough points, and his checks on points are legitimate.
